% ---------------------------------------------------------------------------
% Metodología del VIOG — sección para el paper.
% Todas las cifras salen de data/processed/viog_colombia.csv
% (129 obs, 1994Q1–2026Q1; C-F de DOS COLAS, cf_one_sided=False desde 2026-08-05).
% Verificado: la réplica de los ponderadores reproduce gap_viog y gap_inv_viog
% publicados con error máximo de 4e-15.
% ---------------------------------------------------------------------------
\documentclass[11pt]{article}
\usepackage[spanish]{babel}
\usepackage[utf8]{inputenc}
\usepackage[T1]{fontenc}
\usepackage{booktabs}
\usepackage{amsmath}
\usepackage{geometry}
\geometry{margin=2.5cm}

% Bibliografía: natbib + apalike (disponible en cualquier TeX Live).
% Requiere referencias.bib en la misma carpeta y correr:  pdflatex → bibtex → pdflatex ×2
\usepackage{natbib}
\bibliographystyle{apalike}
% Alternativa recomendada si tienes MacTeX completo: apacite respeta el español
% (escribe "y" en vez de "and" y "et al." correctamente). Comenta las dos
% líneas de arriba y descomenta estas tres:
%   \usepackage[natbibapa]{apacite}
%   \bibliographystyle{apacite}
%   \renewcommand{\bibname}{Referencias}

\title{VIOG: brecha del producto por combinación ponderada de filtros}
\author{}
\date{}

\begin{document}
\maketitle

El VIOG combina cinco extracciones alternativas de la tendencia del PIB colombiano
---Baxter-King \citep{baxter_king_1999}, Christiano-Fitzgerald
\citep{christiano_fitzgerald_2003}, Butterworth, Hodrick-Prescott potenciado
\citep{phillips_shi_2021} y un modelo de componentes no observados estimado por filtro de
Kalman \citep{harvey_1989}--- aplicadas al logaritmo del PIB trimestral
desestacionalizado a precios constantes de 2015 (1994Q1--2026Q1, $n=129$; serie empalmada
por el autor a partir de las bases de Cuentas Nacionales Trimestrales del DANE, que en su
versión vigente arranca en 2005). Cada filtro $i$ entrega una brecha
$g_{i,t}=\ln Y_t-\ln Y^{*}_{i,t}$ y un error absoluto acumulado
$e_{i,t}=N^{-1}\sum_{s\leq t}\lvert g_{i,s}\rvert$. Con ellos se construyen dos compuestos: el
VIOG pondera cada tendencia proporcionalmente a ese error,
$w_{i,t}=e_{i,t}\big/\sum_j e_{j,t}$, de modo que el filtro que atribuye más movimiento al
ciclo pesa más; el 1/VIOG invierte el criterio,
$\tilde w_{i,t}=e_{i,t}^{-1}\big/\sum_j e_{j,t}^{-1}$, y premia al filtro más estable. En
ambos casos el potencial es $\ln Y^{*}_t=\sum_i w_{i,t}\ln Y^{*}_{i,t}$ y la brecha final
$\ln Y_t-\ln Y^{*}_t$. La elección entre uno y otro es poco consecuente en la práctica: las
dos series correlacionan 0,99 y sus desviaciones estándar difieren en 0,1 pp, lo que sugiere
que el resultado descansa en el conjunto de filtros y no en el esquema de ponderación.

Los filtros difieren sobre todo en el tratamiento de los extremos de la muestra, que es
justamente donde la brecha importa para política. Baxter-King, simétrico con $K=12$ rezagos y
adelantos, pierde doce trimestres en cada extremo: no aporta información desde 2023Q2, su
ponderador es exactamente cero en ese tramo y los cuatro filtros restantes se renormalizan.
El Kalman descarta cuatro trimestres de calentamiento y se estima por máxima verosimilitud
sobre $\ln Y$ con tendencia local lineal, ciclo estocástico amortiguado y período acotado
entre 6 y 32 trimestres, con valores iniciales anclados fuera del máximo degenerado de
varianza cíclica nula. Christiano-Fitzgerald se emplea en su versión de dos colas, que en cada
$t$ usa toda la muestra. La elección no lo convierte en filtro de referencia ---ninguno lo es:
los cinco entran al compuesto en pie de igualdad--- sino que lo pone en el mismo terreno que
los otros cuatro, todos ellos suavizados o de dos colas. El costo es que la brecha se revisa
retroactivamente con cada nueva observación \citep{orphanides_vannorden_2002}; la variante
causal, que solo usa $y_1,\dots,y_t$, está disponible como alternativa y sitúa la caída de
2020Q2 en $-2{,}0$ pp en lugar de $-8{,}3$.
El BHP itera tres veces el filtro HP \citep[$\lambda=1600$;][]{hodrick_prescott_1997} sobre el residuo, en vez del criterio de
parada basado en datos de \citet{phillips_shi_2021}, y el Butterworth es un pasa-bajo de
orden 8 con frecuencia de corte
$\omega_c=1/16$ aplicado en fase cero.

El Cuadro~\ref{tab:filtros} resume especificación, cobertura y desempeño. La dispersión de las
brechas ordena los filtros con nitidez: el Kalman atribuye al ciclo cerca de 50\,\% más
variabilidad que Christiano-Fitzgerald (2,86 frente a 1,92 pp de desviación estándar) y, por
construcción del ponderador, termina con el mayor peso del compuesto VIOG (32,6\,\% en 2026Q1)
y con el menor bajo el criterio inverso (18,7\,\%). La divergencia se concentra en 2020Q2: los
dos filtros de banda propiamente dichos ---Baxter-King y Christiano-Fitzgerald--- registran
caídas de 9,3 y 8,3 pp, menos de la mitad de las de Butterworth (17,6), BHP (17,3) y Kalman
(19,3). No es un error de estimación sino una consecuencia de la banda: el colapso y rebote de la
pandemia duró dos o tres trimestres, un \emph{periodo} por debajo del piso de seis y por tanto
una \emph{frecuencia} superior a la que el filtro deja pasar al ciclo. Todo lo que queda fuera
de la banda ---tanto lo más lento como lo más rápido--- cae al residuo que rotulamos tendencia,
de modo que el ``potencial'' de Baxter-King y Christiano-Fitzgerald retrocede 12,9 y 13,0\,\%
en 2020Q2 y rebota cerca de 9,6\,\% en 2020Q3: la desviación estándar de la variación
trimestral de esas dos tendencias es 1,99 y 1,87 pp, frente a 0,44 y 0,40 pp en Butterworth y
BHP y 2,20 pp en el PIB observado. Es una advertencia de lectura más que un defecto ---fuera
de episodios de esa intensidad los cinco filtros coinciden--- y el compuesto sitúa la brecha de
2020Q2 en $-14{,}8$ pp.

\begin{table}[htbp]
\centering
\caption{Filtros del VIOG: especificación, cobertura y desempeño (Colombia, 1994Q1--2026Q1)}
\label{tab:filtros}
\footnotesize
\setlength{\tabcolsep}{3.8pt}
\begin{tabular}{llccrrrr}
\toprule
\textbf{Filtro} & \textbf{Tipo} & \textbf{Parámetro clave} & \textbf{Cobertura} & \textbf{Obs.} &
\textbf{D.E.} & \textbf{2020Q2} & \textbf{Peso} \\
 & & & & & (pp) & (pp) & 2026Q1 \\
\midrule
Baxter-King           & Banda, simétrico  & 6--32 trim., $K{=}12$    & 1997Q1--2023Q1 & 105 & 2,33 & $-9{,}28$  & 0,0\,\%  \\
Christiano-Fitzgerald & Banda, dos colas  & 6--32 trim.              & 1994Q1--2026Q1 & 129 & 1,92 & $-8{,}26$  & 22,3\,\% \\
Butterworth           & Pasa-bajo, fase 0 & orden 8, $\omega_c=1/16$ & 1994Q1--2026Q1 & 129 & 2,30 & $-17{,}62$ & 23,4\,\% \\
HP potenciado (BHP)   & HP iterado        & $\lambda=1600$, 3 iter.  & 1994Q1--2026Q1 & 129 & 2,21 & $-17{,}25$ & 21,7\,\% \\
Kalman / UCM          & Estado-espacio    & 6--32 trim., amort.      & 1995Q1--2026Q1 & 125 & 2,86 & $-19{,}31$ & 32,6\,\% \\
\midrule
\textbf{VIOG}         & Compuesto & $w_{i,t}\propto e_{i,t}$      & 1994Q1--2026Q1 & 129 & 2,19 & $-14{,}78$ & --- \\
\textbf{1/VIOG}       & Compuesto & $w_{i,t}\propto e_{i,t}^{-1}$ & 1994Q1--2026Q1 & 129 & 2,09 & $-14{,}00$ & --- \\
\bottomrule
\end{tabular}

\vspace{0.4em}
\begin{minipage}{0.95\textwidth}
\footnotesize
\textit{Nota:} brechas en puntos porcentuales del PIB ($100\times(\ln Y-\ln Y^{*})$). D.E. es la
desviación estándar de la brecha sobre las observaciones disponibles de cada filtro. El peso de
Baxter-King es cero desde 2023Q2 por el recorte de $K=12$ trimestres en el extremo derecho.
Correlación de cada brecha con el compuesto VIOG: BHP 0,96; Kalman 0,95; Butterworth 0,94;
Baxter-King 0,93; Christiano-Fitzgerald 0,88.
\end{minipage}
\end{table}

\renewcommand{\refname}{Referencias}
\small
\bibliography{referencias}

\end{document}
